% Осень26!B12; нумерация внутри столбца без пустых строк.
% Источник: лекция 1, §14 (строка 378)
% Источник: лекция 2, §2 (строка 88)
% Источник: лекция 2, §3 (строка 113)
% Источник: лекция 2, §4 (строка 131)
\definitionticket{09}{Полукольцо, алгебра, сигма-алгебра}
\label{definition:09}

\begin{definition}
Система множеств $\mathcal P\subset2^X$ называется \textbf{полукольцом}, если:
\begin{enumerate}
  \item $\varnothing\in\mathcal P$;
  \item $A,B\in\mathcal P$ влечёт $A\cap B\in\mathcal P$;
  \item для любых $A,B\in\mathcal P$ разность $A\setminus B$ представима в виде конечного дизъюнктного объединения элементов $\mathcal P$.
\end{enumerate}
\end{definition}

\begin{definition}
Система $\mathcal A\subset2^X$ называется \textbf{алгеброй множеств}, если $X\in\mathcal A$ и для любых $A,B\in\mathcal A$ выполнено $A\setminus B\in\mathcal A$.
\end{definition}

\begin{definition}
Алгебра $\mathcal A$ называется \textbf{$\sigma$-алгеброй}, если она замкнута относительно счётных объединений:
\[
  A_n\in\mathcal A\ (n\in\mathbb N)
  \quad\Longrightarrow\quad
  \bigcup_{n=1}^{\infty}A_n\in\mathcal A.
\]
\end{definition}
